%% ───────────────────────────────────────────────────────────────────
\chapter{La Previdenza Complementare: Il Secondo Pilastro}
%% ───────────────────────────────────────────────────────────────────

\noindent
Il sistema pensionistico pubblico italiano, come dimostrato nei capitoli
precedenti, non è in grado di garantire alle generazioni in regime
contributivo puro un tasso di sostituzione compatibile con il
mantenimento del tenore di vita raggiunto nell'ultima fase lavorativa.
Il gap è strutturale, non temporaneo: non dipende da scelte sbagliate
del singolo lavoratore, ma dalla matematica del sistema NDC in un
contesto di bassa crescita e lunga vita.
In questo scenario, la previdenza complementare cessa di essere un'opzione
per chi ha risorse eccedenti e diventa una necessità strutturale per
chiunque voglia evitare la povertà relativa in età anziana.
Il presente capitolo ha un duplice obiettivo: argomentare questa tesi
con rigore teorico nella prima sezione e fornire al lettore gli strumenti
quantitativi per prendere decisioni informate nelle successive.


%% ───────────────────────────────────────────────────────────────────
\section{Razionale economico del secondo pilastro
         per la generazione NDC}
\label{sec:razionale_2p}
%% ───────────────────────────────────────────────────────────────────

%% ---
\subsection{Il gap previdenziale della Gen Z: tasso di sostituzione
            atteso vs soglia di mantenimento del tenore di vita}
\label{sub:gap_prev}
%% ---

\noindent
La letteratura economica sulla previdenza individua una soglia empirica
al di sotto della quale la pensione non è sufficiente a mantenere il
tenore di vita del periodo lavorativo: il \emph{70--75\% del reddito netto
da lavoro} (sostitution rate target o tasso di sostituzione obiettivo).
Questa soglia non è arbitraria: corrisponde alla riduzione fisiologica
delle spese nella fase pensionistica (minor pendolarismo, nessun
vestiario da lavoro, nessuna contribuzione) al netto delle nuove voci
di costo (sanità, cura della persona, tempo libero) \cite{Bosi2018}.

\medskip
\noindent
Le proiezioni RGS 2024--2025 documentano che le generazioni in regime
contributivo puro — chi ha iniziato a lavorare dopo il 1996 —
si collocheranno strutturalmente al di sotto di questa soglia
\cite{MEFRGS2025}:

\begin{table}[h!]
\centering
\small
\caption{Gap previdenziale per generazione: tasso di sostituzione NDC atteso
         vs soglia obiettivo 70\% (scenario base RGS, carriera continua 40 anni)}
\label{tab:gap_prev}
\begin{tabular}{lcccl}
\toprule
\textbf{Coorte} & \textbf{Anno pensionamento} &
\textbf{Tasso sostituz. atteso} & \textbf{Gap vs 70\%}
& \textbf{Profilo contributivo} \\
\midrule
Baby Boomers    & 2005--2020 & 79,3\%  & +9,3 pp  & Retributivo pieno \\
Gen X (ibrida)  & 2025--2035 & 72,1\%  & +2,1 pp  & Misto (pro-rata) \\
Millennials     & 2040--2050 & 64,8\%  & --5,2 pp & Contributivo puro \\
Gen Z           & 2055--2065 & 58--63\%& --7--12 pp& Contributivo puro + discontinuo \\
\bottomrule
\multicolumn{5}{l}{\small Fonti: \cite{MEFRGS2025, SlideCorso2025,
CensisConfcooperative2026}.}
\end{tabular}
\end{table}

\noindent
Il dato più rilevante non è il tasso di sostituzione atteso in isolamento,
ma il \emph{gap da colmare}: per la Gen Z con carriera continua, il gap
è di 7--12 punti percentuali rispetto alla soglia minima;
per chi ha carriera discontinua (il profilo più diffuso tra i lavoratori
nati dopo il 1990, come mostrato nel Capitolo~3), il gap supera i
15--20 punti.
Un lavoratore Gen Z con reddito netto da lavoro di 1.800 euro/mese
(livello mediano italiano per i 35--45enni, 2024) che vada in pensione
con un tasso di sostituzione del 60\% riceve 1.080 euro/mese:
il divario di 720 euro mensili rispetto alla soglia del 70\%
deve essere colmato da qualche altra fonte di reddito —
o non verrà colmato affatto \cite{MEFRGS2025, Bosi2018}.

%% ---
\subsection{L'effetto sostituzione della ricchezza nel PAYG:
            perché il lavoratore NDC deve risparmiare
            attivamente più delle generazioni precedenti}
\label{sub:effetto_sostituzione}
%% ---

\noindent
La teoria economica del ciclo vitale identifica un meccanismo —
l'\emph{effetto sostituzione della ricchezza} — per cui la previdenza
pubblica riduce il risparmio privato volontario: l'individuo che
si aspetta di ricevere una pensione adeguata risparmia meno nell'immediato,
trattando i contributi obbligatori come parte del proprio piano
di risparmio complessivo \cite[slide~19]{SlideCorso2025}.

\medskip
\noindent
Nel sistema retributivo, questo meccanismo produceva un risultato
accettabile: l'individuo risparmiava meno perché la pensione attesa
(80\% dell'ultima retribuzione) rendeva superfluo un risparmio
complementare.
Nel sistema NDC con tasso di sostituzione del 60\%, la situazione
è opposta: il risparmio privato non viene spiazzato abbastanza.
L'individuo continua a versare il 33\% della retribuzione lorda
in contributi obbligatori, sviluppando la percezione psicologica di
``pagare per la propria pensione'', ma il ritorno atteso di quei
contributi è strutturalmente insufficiente \cite{Bosi2018, SlideCorso2025}.

\medskip
\noindent
L'effetto sostituzione nel NDC crea quindi un paradosso comportamentale:
\begin{enumerate}
    \item L'aliquota contributiva del 33\% crea l'\emph{illusione}
          di un risparmio previdenziale adeguato.
    \item In realtà, il rendimento implicito di quei contributi
          (approssimativamente $g \approx 0{,}5\%$ reale) è inferiore
          al rendimento di mercato, e la pensione attesa è insufficiente.
    \item Il lavoratore che non si informa e non agisce in modo
          integrativo arriverà al pensionamento con uno stile di vita
          incompatibile con il reddito pensionistico effettivo.
\end{enumerate}

\medskip
\noindent
Il corollario normativo è preciso: il lavoratore Gen Z che non aderisce
alla previdenza complementare non sta semplicemente rinunciando a un
beneficio aggiuntivo; sta scommettendo che la propria pensione NDC,
calcolata su una speranza di vita di 87--90 anni, sarà sufficiente a
mantenere il suo tenore di vita.
Una scommessa che, con probabilità elevata, perderà \cite{Fornero2018}.

%% ---
\subsection{Fondi negoziali, aperti e PIP:
            struttura, costi medi e rendimenti storici
            a confronto (dati COVIP 2024)}
\label{sub:tipologie}
%% ---

\noindent
Il sistema italiano di previdenza complementare si articola, ai sensi
del D.Lgs.\ 252/2005, in tre categorie principali, con strutture
di costi, rendimenti e accessibilità significativamente differenti
\cite{COVIP2025}.

\begin{table}[h!]
\centering
\small
\caption{Tipologie di previdenza complementare in Italia: confronto
         strutturale su cinque dimensioni (dati COVIP 2024)}
\label{tab:tipologie_fondi}
\begin{tabular}{p{2.8cm} p{3.5cm} p{3.5cm} p{3.5cm}}
\toprule
\textbf{Dimensione}
& \textbf{Fondo Negoziale}
& \textbf{Fondo Aperto}
& \textbf{PIP (Piano Individuale Pensionistico)} \\
\midrule
\textbf{Natura}
& Chiuso, istituito da contrattazione collettiva
& Aperto a tutti, istituito da banche/SGR
& Prodotto assicurativo (ramo vita III o I) \\
\addlinespace
\textbf{Accessibilità}
& Solo per dipendenti del settore di riferimento
& Chiunque (dipendenti, autonomi, liberi professionisti)
& Chiunque; spesso proposto da reti bancarie e agenti \\
\addlinespace
\textbf{ISC (Indicatore Sintetico Costi) a 10 anni}
& \textbf{0,49\%}
& 1,35\%
& 2,17\% \\
\addlinespace
\textbf{Rendimento linea bilanciata 2024}
& 6,4\%
& 6,6\%
& 7,0\% \\
\addlinespace
\textbf{Rendimento linea azionaria 2024}
& 10,4\%
& 10,4\%
& 13,0\% \\
\addlinespace
\textbf{Rendimento medio annuo 2015--2024}
& 2,2\%
& 2,4\%
& 2,9\% \\
\addlinespace
\textbf{Impatto costo su montante a 35 anni}
& Riduzione $\approx$5\%
& Riduzione $\approx$22\%
& Riduzione $\approx$40\% \\
\addlinespace
\textbf{Contributo datore di lavoro}
& Sì (obbligatorio se lavoratore aderisce)
& No (solo TFR)
& No \\
\addlinespace
\textbf{Iscritti a fine 2024}
& 4.109.000 (+5,5\%)
& 1.800.000 (+7,0\%)
& 3.600.000 (+2,5\%) \\
\bottomrule
\multicolumn{4}{l}{\small Fonti: \cite{COVIP2025, web:65, web:80}.}
\end{tabular}
\end{table}

\noindent
Il dato più critico nella Tabella~\ref{tab:tipologie_fondi} è l'impatto
dei costi sul montante a lungo termine: su un orizzonte di 35 anni,
una differenza dell'1\% annuo di costo di gestione produce una
riduzione del montante finale del 18--20\% secondo le simulazioni COVIP
\cite{web:80}.
Per un montante atteso di 200.000 euro, significa una perdita di
36.000--40.000 euro attribuibile esclusivamente alla scelta del
veicolo più costoso.
La conclusione operativa per il lavoratore Gen Z è diretta:
se ha accesso a un fondo negoziale di categoria, non ci sono argomenti
razionali per scegliere un PIP — che ha un ISC di 2,17\% contro lo
0,49\% del negoziale — a meno che non abbia caratteristiche specifiche
irripetibili in altri veicoli.

\medskip
\noindent
A fine 2024, il totale degli iscritti alle forme di previdenza
complementare ha superato per la prima volta i \textbf{10 milioni}
(9.952.654 posizioni, equivalenti a circa 8,7 milioni di individui),
con risorse accumulate pari a \textbf{243,4 miliardi di euro}
(10,8\% del PIL) \cite{web:65, COVIP2025}.
Ma la percentuale di adesione rispetto alla forza lavoro si attesta
al 38,3\% — con solo il 27,6\% degli iscritti che ha effettivamente
versato contributi nel 2024 \cite{web:85}.
Il 62\% dei lavoratori italiani non aderisce ad alcuna forma di
previdenza complementare: è la cifra che sintetizza il fallimento
informativo che questo capitolo intende affrontare \cite{web:76}.


%% ───────────────────────────────────────────────────────────────────
\section{TFR in azienda vs fondo pensione: analisi di convenienza}
\label{sec:tfr_fondo}
%% ───────────────────────────────────────────────────────────────────

%% ---
\subsection{Il TFR come ``risparmio forzato silente'':
            meccanismo, rivalutazione e costo-opportunità
            del non investirlo}
\label{sub:tfr_meccanismo}
%% ---

\noindent
Il Trattamento di Fine Rapporto (TFR) è uno degli istituti più peculiari
del diritto del lavoro italiano: una forma di retribuzione differita,
accantonata annualmente in misura pari al 6,91\% della retribuzione
lorda e liquidata al momento della cessazione del rapporto di lavoro
(per qualsiasi causa: pensionamento, dimissioni, licenziamento).
Per un lavoratore con stipendio lordo di 25.000 euro/anno, il TFR
maturato annualmente è circa \textbf{1.727 euro}: una somma non irrisoria,
che in 40 anni di carriera forma un capitale significativo.

\medskip
\noindent
Il TFR non è quindi ``denaro del datore di lavoro'': è retribuzione
del lavoratore, accantonata obbligatoriamente e rivalutata ogni anno
secondo la formula:
\begin{equation}
    \text{Rivalutazione TFR} = 1{,}5\% + 0{,}75 \times \pi_{\text{ISTAT}}
    \label{eq:rivalutazione_tfr}
\end{equation}
dove $\pi_{\text{ISTAT}}$ è il tasso di inflazione annua rilevato dall'ISTAT.
Storicamente, con inflazione media del 2\%, il tasso di rivalutazione
netto del TFR si è attestato intorno al 2,5\% \cite{web:86, web:65}.

\medskip
\noindent
Il TFR lasciato in azienda presenta quattro svantaggi strutturali:
\begin{enumerate}
    \item \emph{Rendimento basso}: 2,4\% medio annuo netto nel periodo
          2015--2024, contro il 2,2--2,9\% dei fondi pensione nello
          stesso periodo — ma con profilo di rischio molto diverso:
          in anni di alta inflazione (2022: +8,1\%), il TFR in azienda
          sovraperforma; in anni di bassa inflazione e mercati azionari
          positivi (2023, 2024), il fondo pensione è nettamente superiore
          \cite{web:83, web:66}.
    \item \emph{Tassazione svantaggiosa al momento della liquidazione}:
          il TFR è tassato con l'\emph{aliquota media IRPEF degli ultimi
          5 anni} del lavoratore, spesso pari al 23--38\%, con
          possibilità di tassazione separata ma comunque non inferiore
          al 23\% per i redditi mediani.
    \item \emph{Rischio d'impresa}: il TFR depositato presso aziende
          con meno di 50 dipendenti (la grande maggioranza in Italia)
          non è garantito dal Fondo di Garanzia INPS in caso di
          insolvenza, esponendo il lavoratore a un rischio di credito
          non remunerato.
    \item \emph{Nessun contributo del datore di lavoro}: chi lascia
          il TFR in azienda perde il contributo aggiuntivo del datore
          che i fondi negoziali prevedono obbligatoriamente per i
          lavoratori aderenti.
\end{enumerate}

%% ---
\subsection{Il vantaggio fiscale del fondo pensione: tassazione
            15\%--9\% vs aliquota marginale IRPEF,
            deducibilità e contributo del datore}
\label{sub:vantaggio_fiscale}
%% ---

\noindent
Il fondo pensione presenta un trattamento fiscale privilegiato
articolato su tre livelli, ciascuno dei quali genera un vantaggio
quantificabile \cite{web:78, web:84, web:93}.

\medskip
\noindent
\textbf{Livello 1 — Deducibilità dei contributi volontari.}
I contributi versati dal lavoratore (non il TFR, che segue un
regime separato) sono deducibili dall'imponibile IRPEF fino al
limite massimo di \textbf{5.300 euro/anno} (aggiornato dalla
Legge di Bilancio 2026, rispetto ai precedenti 5.164,57 euro)
\cite{web:93}.
In termini concreti: un lavoratore con reddito imponibile di 30.000 euro
(aliquota marginale IRPEF al 35\%) che versi 2.000 euro/anno al fondo
risparmia 700 euro di IRPEF — un rendimento immediato del 35\%
sul versamento, prima ancora di considerare il rendimento finanziario.

\medskip
\noindent
\textbf{Livello 2 — Tassazione agevolata in fase di erogazione.}
Le prestazioni erogate dal fondo pensione sono tassate con un'aliquota
del \textbf{15\%}, che si riduce dello 0,30\% per ogni anno di
partecipazione oltre il quindicesimo, fino al \textbf{9\% minimo}
raggiungibile dopo 35 anni di partecipazione.
Il confronto con la tassazione del TFR (23--38\%) rende immediatamente
percepibile il vantaggio: uno stesso capitale accumulato di 100.000 euro
lascia nelle tasche del pensionato 91.000 euro se liquidato dal fondo
(dopo 35 anni di partecipazione), contro 62.000--77.000 euro se liquidato
come TFR in azienda \cite{web:78}.

\medskip
\noindent
\textbf{Livello 3 — Il contributo del datore di lavoro: denaro gratuito.}
Per i lavoratori dipendenti che aderiscono al fondo negoziale
di categoria, il contratto collettivo prevede che il datore di lavoro
versi una quota aggiuntiva — tipicamente l'1--1,5\% della retribuzione lorda —
che viene accreditata sulla posizione del lavoratore
\emph{soltanto se il lavoratore stesso aderisce e versa il minimo contrattuale}.
Per un lavoratore con stipendio di 25.000 euro, il contributo del datore
vale circa 250--375 euro/anno: se non si aderisce, questo denaro
viene semplicemente non percepito.
In 40 anni, capitalizzato al 4\%, il contributo del datore da solo
vale 24.000--36.000 euro di montante \cite{COVIP2025}.
Rinunciare al fondo negoziale quando si è dipendenti equivale a
rifiutare parte della propria retribuzione contrattuale.

%% ---
\subsection{Simulazione numerica: montante a 40 anni di carriera --
            TFR in azienda vs fondo negoziale, euro per euro
            e scenario con/senza contributo datore}
\label{sub:simulazione}
%% ---

\noindent
La seguente simulazione confronta in modo sistematico i due scenari
per il lavoratore tipo della tesi (Marco, 24 anni, stipendio 25.000
euro/anno, 40 anni di carriera, pensionamento a 64 anni).

\medskip
\noindent\textbf{Parametri della simulazione:}

\begin{table}[h!]
\centering
\small
\caption{Parametri di simulazione TFR in azienda vs fondo negoziale}
\label{tab:parametri_sim}
\begin{tabular}{lll}
\toprule
\textbf{Parametro} & \textbf{Caso A -- TFR in azienda}
                   & \textbf{Caso B -- Fondo negoziale} \\
\midrule
Accantonamento annuo & €1.727 (TFR 6,91\% × €25.000) & €2.727 \\
\quad di cui: TFR conferito & €1.727 & €1.727 \\
\quad di cui: contributo datore & --- & €500 \\
\quad di cui: contributo lavoratore & --- & €500 \\
Tasso di rendimento/rivalutazione & 2,5\% (formula TFR) & 4,0\% (bilanciato) \\
Costo annuo & 0\% & 0,49\% (ISC negoziale) \\
Aliquota fiscale alla liquidazione & 27\% (IRPEF media) & 11\% (15\%--9\%) \\
Durata & 40 anni & 40 anni \\
\bottomrule
\end{tabular}
\end{table}

\medskip
\noindent\textbf{Risultati della simulazione:}

\begin{table}[h!]
\centering
\small
\caption{Confronto montante netto: TFR in azienda vs fondo negoziale (40 anni)}
\label{tab:risultati_sim}
\begin{tabular}{lcc}
\toprule
\textbf{Scenario} & \textbf{Montante lordo} & \textbf{Montante netto} \\
\midrule
A -- TFR in azienda
    & €116.438 & \textbf{€85.000} \\
B1 -- Fondo negoziale (solo TFR, no contributo datore)
    & €164.420 & \textbf{€146.334} \\
B2 -- Fondo negoziale (TFR + datore + lavoratore)
    & €259.182 & \textbf{€230.672} \\
\midrule
\multicolumn{2}{l}{\textbf{Vantaggio B1 vs A (stesso esborso)}}
    & \textbf{+€61.334 (+72\%)} \\
\multicolumn{2}{l}{\textbf{Vantaggio B2 vs A (contributo datore+lavoratore)}}
    & \textbf{+€145.672 (+171\%)} \\
\bottomrule
\multicolumn{3}{l}{\small Elaborazione propria su parametri COVIP 2024 \cite{COVIP2025, web:65}.}
\end{tabular}
\end{table}

\begin{figure}[h!]
\centering
% placeholder for chart:96
\caption{Accumulo netto a 40 anni: TFR in azienda vs fondo negoziale
         (lavoratore tipo €25.000/anno)}
\label{fig:tfr_fondo}
\small\emph{Nota: Elaborazione su dati COVIP 2024. Caso A: TFR rivalutato al 2,5\%,
tassazione IRPEF 27\%. Caso B: TFR+contributo datore (€500)+lavoratore (€500),
rendimento bilanciato 4\%, ISC 0,49\%, tassazione 11\%.}
\end{figure}

\noindent
Il confronto tra B1 e A è particolarmente rilevante perché misura
il vantaggio \emph{a parità di flussi di cassa del lavoratore}:
anche senza versare un euro aggiuntivo oltre al TFR già maturato,
la sola destinazione al fondo negoziale produce un montante netto
del 72\% superiore rispetto al TFR lasciato in azienda.
La differenza è interamente spiegabile dal rendimento superiore (4\%
vs 2,5\%) e dalla tassazione più favorevole (11\% vs 27\%) \cite{COVIP2025}.

%% ---
\subsection{Il costo del ritardo: quanto vale aspettare 5 o 10 anni
            prima di aderire (formula $M = C \cdot (1+r)^n$)}
\label{sub:costo_ritardo}
%% ---

\noindent
L'interesse composto è il principio finanziario più potente e più
sistematicamente ignorato nel comportamento previdenziale italiano.
La formula del montante di una rendita a rate costanti posticipate è:
\begin{equation}
    M = C \cdot \frac{(1+r)^n - 1}{r}
    \label{eq:montante_rendita}
\end{equation}
dove $C$ è il versamento periodico, $r$ il tasso di rendimento e
$n$ il numero di periodi.
La proprietà fondamentale della formula \eqref{eq:montante_rendita}
è la sua \emph{convessità in $n$}: all'aumentare del numero di periodi,
il montante cresce in modo più che proporzionale, perché gli interessi
degli ultimi anni maturano su una base sempre più larga.
I primi anni di contribuzione valgono più degli ultimi, non perché
i contributi siano più grandi, ma perché il tempo di capitalizzazione
è massimo.

\medskip
\noindent
La Tabella \ref{tab:costo_ritardo} quantifica il costo del ritardo
per il lavoratore tipo (versamento di 200 euro/mese, rendimento 4\%):

\begin{table}[h!]
\centering
\small
\caption{Costo del ritardo: montante a 67 anni per diversa età di inizio
         versamento (€200/mese, r=4\% annuo)}
\label{tab:costo_ritardo}
\begin{tabular}{ccccc}
\toprule
\textbf{Età di inizio} & \textbf{Anni di contribuzione} &
\textbf{Capitale versato} & \textbf{Montante a 67 anni} &
\textbf{Perdita vs inizio a 22} \\
\midrule
22 & 45 & €108.000 & \textbf{€295.758} & --- \\
27 & 40 & €96.000  & \textbf{€232.213} & --€63.545 (--21\%) \\
32 & 35 & €84.000  & \textbf{€179.983} & --€115.775 (--39\%) \\
37 & 30 & €72.000  & \textbf{€137.054} & --€158.704 (--54\%) \\
42 & 25 & €60.000  & \textbf{€101.770} & --€193.988 (--66\%) \\
\bottomrule
\multicolumn{5}{l}{\small Elaborazione propria. Fonte: formula \eqref{eq:montante_rendita}.}
\end{tabular}
\end{table}

\noindent
La Tabella \ref{tab:costo_ritardo} rivela un dato sorprendente:
chi inizia a 42 anni versa in totale 60.000 euro di contributi
in meno rispetto a chi inizia a 22, ma perde quasi 194.000 euro
di montante finale — un moltiplicatore della perdita pari a 3,2 volte
il capitale non versato.
Non è il capitale il problema: è il tempo di capitalizzazione.
Ogni anno di ritardo nell'iscrizione al fondo pensione vale in media
\textbf{circa 10.000--15.000 euro di montante finale perduto},
un costo silenzioso e irreversibile che nessuna contribuzione futura
può recuperare integralmente.

%% ---
\subsection{Solo il 38,3\% dei lavoratori aderisce alla previdenza
            complementare: analisi delle cause e implicazioni di policy}
\label{sub:bassa_adesione}
%% ---

\noindent
A fronte di vantaggi così documentati, il tasso di adesione alla
previdenza complementare in Italia appare incomprensibile:
il 62\% della forza lavoro non ha alcuna forma di secondo pilastro attivo,
e di chi è formalmente iscritto, ben il 27,7\% non ha versato contributi
nel 2024 \cite{web:76, web:85, COVIP2025}.
Quattro cause strutturali spiegano questa persistente bassa adesione.

\medskip
\noindent
\textbf{Causa 1 — Vincolo di liquidità.}
Il versamento al fondo pensione, anche minimo, richiede un flusso di
cassa corrente disponibile.
Per un lavoratore con reddito basso e spese correnti elevate (affitto,
figlio, mutuo), destinare anche solo 50--100 euro/mese al fondo
può essere oggettivamente impossibile nel breve periodo, indipendentemente
dalla sua razionalità finanziaria.
Questo spiega la correlazione documentata tra reddito e adesione
al secondo pilastro: chi guadagna di più aderisce di più, amplificando
la disuguaglianza previdenziale già presente nel primo pilastro
\cite{Bosi2018, COVIP2025}.

\medskip
\noindent
\textbf{Causa 2 — Inerzia comportamentale e status quo bias.}
La psicologia economica (Kahneman, Thaler) ha documentato
che le decisioni che richiedono un'azione attiva di iscrizione
(sistema \emph{opt-in}) sono sistematicamente procrastinate
rispetto a decisioni reversibili che non richiedono azione
(sistema \emph{opt-out}).
Il silenzio-assenso introdotto dal D.Lgs.\ 252/2005 per la
destinazione del TFR ai fondi pensione si applica solo alle
nuove assunzioni e solo per la quota TFR: non copre il versamento
volontario e non si estende ai lavoratori già assunti al momento
della sua introduzione \cite{web:80, Lavoce2024a}.

\medskip
\noindent
\textbf{Causa 3 — Ignoranza strutturale.}
L'indagine COVIP 2023 sulla cultura previdenziale documenta che
oltre il 50\% dei lavoratori italiani non conosce l'esistenza
del fondo negoziale di categoria di riferimento,
e il 70\% non sa qual è l'aliquota contributiva minima per
ottenere il contributo del datore \cite{COVIP2025}.
Non si tratta di irrazionalità: si tratta di un fallimento informativo
che il sistema non corregge perché nessun agente di mercato ha
interesse a informare il lavoratore circa strumenti a basso costo
gestiti non dal settore privato \cite{Lavoce2024a}.

\medskip
\noindent
\textbf{Causa 4 — Diffidenza istituzionale.}
I fallimenti di singoli fondi pensione (anni '90), la memoria
del blocco dei riscatti durante la crisi finanziaria del 2008 e la
percezione di un sistema opaco e poco flessibile nelle modalità di
accesso (il capitale è vincolato fino al pensionamento, con eccezioni
limitate) creano una diffidenza razionale verso strumenti che
immobilizzano liquidità per decenni \cite{Fornero2018}.


%% ───────────────────────────────────────────────────────────────────
\section{L'interesse composto e il tempo come asset}
\label{sec:interesse_composto}
%% ───────────────────────────────────────────────────────────────────

%% ---
\subsection{La matematica del tempo: $M = C \cdot (1+r)^n$
            e il valore dei primi 10 anni di carriera}
\label{sub:matematica_tempo}
%% ---

\noindent
La formula del montante di un capitale iniziale
\begin{equation}
    M = C_0 \cdot (1+r)^n
    \label{eq:montante_cap}
\end{equation}
è matematicamente semplice ma psicologicamente controintuitiva.
Il problema non è la formula: è il fatto che l'impatto del fattore
$(1+r)^n$ è percepito come lineare nell'immaginario comune, mentre
nella realtà è esponenziale.
Un euro investito a 22 anni al 4\% vale, a 67 anni:
\[
    M = 1 \cdot (1{,}04)^{45} = 5{,}84 \text{ euro}
\]
Lo stesso euro investito a 42 anni vale a 67 anni solo:
\[
    M = 1 \cdot (1{,}04)^{25} = 2{,}67 \text{ euro}
\]
La differenza di 20 anni di tempo (non di contributo aggiuntivo)
moltiplica il rendimento finale per 2,2 volte.
I \emph{primi 10 anni di carriera} hanno quindi un valore previdenziale
sproporzionatamente alto: sono gli anni in cui il tempo di
capitalizzazione è massimo e il costo-opportunità del ritardo
è massimo.

\medskip
\noindent
Per rendere tangibile questo principio, si confronti il montante
di chi versa 200 euro/mese per 40 anni (da 22 a 62, poi smette)
con quello di chi versa gli stessi 200 euro/mese per 30 anni
(da 32 a 62) al tasso del 4\%:
\begin{align}
    M_{40 \text{ anni}} &= 200 \cdot
    \frac{(1+r_m)^{480}-1}{r_m} \approx \text{€}232.213
    \label{eq:m40} \\
    M_{30 \text{ anni}} &= 200 \cdot
    \frac{(1+r_m)^{360}-1}{r_m} \approx \text{€}139.316
    \label{eq:m30}
\end{align}
dove $r_m = (1{,}04)^{1/12}-1$ è il tasso mensile equivalente.
I \textbf{10 anni di differenza} (da 22 a 32) producono una differenza
di montante di quasi \textbf{93.000 euro}: l'equivalente di quasi
4 anni di stipendio del lavoratore tipo, generati non da versamenti
aggiuntivi (i 200 euro/mese sono uguali in entrambi i casi)
ma dalla sola forza dell'interesse composto nei primi anni di carriera.

%% ---
\subsection{Il grafico del vantaggio temporale: montante a 67 anni
            in funzione dell'età di inizio contribuzione}
\label{sub:grafico_tempo}
%% ---

\begin{figure}[h!]
\centering
% placeholder for chart:97
\caption{Il costo del ritardo: montante a 67 anni per €200/mese al 4\%
         in funzione dell'età di inizio contribuzione}
\label{fig:vantaggio_tempo}
\small\emph{Nota: Rendimento annuo 4\% (bilanciato fondo negoziale), versamento €200/mese.
L'asse verticale mostra il montante lordo nozionale;
il montante netto risulta circa il 15-20\% inferiore per effetto
della tassazione al pensionamento.}
\end{figure}

\noindent
La Figura~\ref{fig:vantaggio_tempo} mostra visivamente la struttura
del problema: il montante non decresce linearmente con il ritardo,
ma in modo accelerato.
Passare da 22 a 27 anni (5 anni di ritardo) costa 63.545 euro
di montante finale; passare da 37 a 42 anni (stessi 5 anni)
costa ``solo'' 35.284 euro.
I primi anni sono i più costosi da perdere: questo è il principio
che dovrebbe guidare qualsiasi politica di incentivo all'adesione
precoce dei giovani lavoratori.

%% ---
\subsection{L'iscrizione automatica (\emph{auto-enrollment})
            come soluzione all'inerzia:
            evidenze da USA, UK e proposta per l'Italia}
\label{sub:autoenrollment}
%% ---

\noindent
La ricerca sperimentale in economia comportamentale (Thaler e Sunstein,
\emph{Nudge}, 2008) ha documentato che la scelta predefinita
(\emph{default}) influenza in modo determinante il comportamento
individuale in presenza di inerzia e complessità decisionale.
Applicato alla previdenza complementare, il principio implica che
passare da un sistema \emph{opt-in} (il lavoratore deve agire attivamente
per iscriversi) a un sistema \emph{opt-out} (il lavoratore è
automaticamente iscritto e deve agire per uscire) produce
un aumento dell'adesione drammatico e persistente nel tempo.

\medskip
\noindent
\textbf{L'evidenza USA (401(k) plan).}
Negli Stati Uniti, l'introduzione dell'auto-enrollment nei piani 401(k)
a partire dagli anni 2000 ha aumentato il tasso di adesione
dei nuovi assunti dal 37--42\% (opt-in) all'83--90\% (opt-out)
in modo stabile nel tempo.
L'esperimento più documentato, condotto da Madrian e Shea (2001)
su una grande azienda manifatturiera USA, ha mostrato un salto
nell'adesione di 49 punti percentuali nel primo mese di impiego
semplicemente modificando il default da ``non iscritto'' a ``iscritto
con contribuzione al 3\% nella linea bilanciata predefinita'' \cite{Lavoce2024a}.

\medskip
\noindent
\textbf{L'evidenza UK (NEST — National Employment Savings Trust).}
Nel 2012, il governo britannico ha introdotto l'\emph{auto-enrolment}
obbligatorio per tutti i datori di lavoro con almeno un dipendente.
Ogni lavoratore tra i 22 e i 66 anni con reddito superiore a
£10.000 viene automaticamente iscritto al fondo aziendale (o al
fondo NEST se l'azienda non ha un fondo proprio) con un contributo
minimo del 5\% del reddito (3\% datore + 2\% lavoratore).
Il tasso di opt-out è rimasto stabilmente intorno al 9\%:
oltre il 91\% dei lavoratori ha mantenuto l'iscrizione automatica.
In 12 anni, il numero di lavoratori con secondo pilastro attivo
è passato da 10 milioni a oltre 22 milioni \cite{Lavoce2024a, Fornero2018}.

\medskip
\noindent
\textbf{La proposta per l'Italia: adesione automatica per i nuovi assunti.}
La Legge di Bilancio 2026 ha recepito, per la prima volta in forma
strutturale, una proposta di meccanismo di adesione automatica per i
\emph{nuovi assunti}: i datori di lavoro sarebbero tenuti a iscrivere
automaticamente i nuovi dipendenti al fondo negoziale di categoria
(o, in assenza, a un fondo aperto indicato dal CCNL), con contribuzione
predefinita alla linea life-cycle adeguata all'età anagrafica
\cite{web:70, web:73}.
Il lavoratore mantiene il diritto di recesso entro 30 giorni senza
penalità.
La misura — se approvata e implementata con la continuità necessaria —
potrebbe aumentare il tasso di adesione degli under-35 dall'attuale
circa 25\% verso il 70--80\% nel medio periodo, avvicinando l'Italia
ai risultati del modello britannico \cite{web:80, COVIP2025}.

\medskip
\noindent
La critica più frequente all'auto-enrollment è che costituisce
una forma di paternalismo incompatibile con la libertà di scelta
individuale.
La risposta teorica è quella di Thaler e Sunstein: l'\emph{opzione
di default} non è mai neutrale.
In un sistema opt-in, il default è ``non iscritto'': è una scelta
che penalizza il lavoratore inerte a vantaggio dei provider
finanziari che hanno interesse a vendere prodotti sostitutivi.
In un sistema opt-out con default calibrato sul fondo negoziale
a basso costo, il default favorisce la scelta che la teoria
economica identifica come ottimale per la grande maggioranza
dei lavoratori.
Non è paternalismo: è architettura della scelta, orientata
verso l'interesse del lavoratore invece che verso i margini
del settore finanziario \cite{Lavoce2024a, Bosi2018}.

\bigskip
\noindent\rule{\textwidth}{0.4pt}

\medskip
\noindent\textbf{\large Nota operativa conclusiva per il lavoratore Gen Z:
                          la lista delle cinque azioni}

\medskip
\noindent
Al termine di questo capitolo, il lettore dispone di tutti gli strumenti
teorici e quantitativi per comprendere il secondo pilastro.
È utile tradurli in una lista operativa concisa — non come consulenza
finanziaria personalizzata, ma come sintesi del messaggio analitico
del capitolo:

\begin{enumerate}
    \item \textbf{Scopri il tuo fondo negoziale di categoria.}
          Esiste per la quasi totalità dei lavoratori dipendenti
          italiani ed è il veicolo con il costo più basso (ISC 0,49\%),
          il contributo del datore di lavoro incluso e la massima
          tutela regolamentare.
    \item \textbf{Versa almeno il minimo contrattuale.}
          Il contributo minimo del lavoratore (tipicamente 1--1,5\%
          della retribuzione) è la soglia che sblocca il contributo
          del datore: sotto questa soglia stai rifiutando retribuzione
          contrattuale.
    \item \textbf{Destina il TFR al fondo pensione.}
          Come mostrato nella Tabella~\ref{tab:risultati_sim}, la sola
          destinazione del TFR (senza versamenti aggiuntivi) produce un
          montante netto del 72\% superiore rispetto al TFR lasciato
          in azienda.
    \item \textbf{Non aspettare.}
          Ogni anno di ritardo nell'iscrizione costa in media
          10.000--15.000 euro di montante finale perduto.
          Non esistono versamenti futuri che compensino completamente
          il tempo perso.
    \item \textbf{Scegli la linea bilanciata o azionaria a 30--40 anni
          dalla pensione.}
          L'orizzonte temporale lungo rende il rischio di mercato
          ampiamente gestibile: una linea garantita o obbligazionaria
          a 30 anni dalla pensione è, paradossalmente, la scelta
          più rischiosa in termini di rendimento atteso insufficiente.
\end{enumerate}